\section{Clasificación de los datos según su tipo}

La clasificación considera forma, función y sensibilidad. Una entidad puede
pertenecer a más de una categoría: una consulta es un dato operativo y también
puede alimentar análisis posteriores.

\begin{table}[ht]
\centering
\small
\begin{tabularx}{\textwidth}{>{\raggedright\arraybackslash}p{0.22\linewidth}Y}
\toprule
Tipo & Ejemplos en la solución \\
\midrule
Estructurados & Catálogos, productos, clientes, proveedores, condiciones comerciales, pedidos, entregas, permisos, consultas y evidencias. Se representan con columnas tipadas, claves y restricciones. \\
Semiestructurados & Metadatos de versiones, parámetros y contenido de resultados estructurados, y detalles de auditoría. Se guardan en JSONB porque su forma puede variar dentro de límites controlados. \\
No estructurados & Archivos Markdown de las versiones documentales. Permanecen fuera de PostgreSQL y se vinculan mediante ruta, tipo MIME, tamaño y checksum. \\
Vectoriales & Embeddings sintéticos de los fragmentos documentales, almacenados como \texttt{VECTOR(32)} y asociados a una versión identificable del modelo. \\
Operativos & Pedidos, líneas, condiciones comerciales, entregas, incidencias, consultas y respuestas. Sostienen la operación o registran una interacción concreta. \\
Analíticos & Totales por categoría, rankings por segmento, métricas de uso y planes de ejecución. En la prueba se calculan bajo demanda y no justifican un Data Warehouse separado. \\
Sensibles o restringidos & Condiciones comerciales, incidencias y documentos legales o de cumplimiento. Su visibilidad depende del perfil y de la clase documental. \\
Auditoría y trazabilidad & Eventos de consulta, acceso, publicación, recuperación y evidencia. Conservan referencias y contexto, pero no duplican el contenido sensible. \\
\bottomrule
\end{tabularx}
\caption{Clasificación de los datos de la solución.}
\end{table}

Los datos analíticos no se guardan por separado porque el volumen y las
consultas del TP no lo justifican. Si el historial de uso o la auditoría
afectaran al núcleo transaccional, podrían replicarse en otra capa sin cambiar
el modelo operativo.
